\documentclass[11pt,a4paper]{article}
% Preâmbulo institucional comum — DEE/IFMA
\usepackage[utf8]{inputenc}
\usepackage[T1]{fontenc}
\usepackage[brazil]{babel}
\usepackage{lmodern}
\usepackage{microtype}
\usepackage{geometry}
\usepackage{xcolor}
\usepackage{graphicx}
\usepackage{booktabs}
\usepackage{tabularx}
\usepackage{array}
\usepackage{enumitem}
\usepackage{hyperref}
\usepackage{lastpage}
\usepackage{fancyhdr}
\usepackage{setspace}
\usepackage{titlesec}
\usepackage{tcolorbox}
\usepackage{ragged2e}

\geometry{a4paper,top=2.2cm,bottom=2.2cm,left=2.5cm,right=2.5cm}
\definecolor{azulpetroleo}{HTML}{0F4C5C}
\definecolor{ambar}{HTML}{C8892D}
\definecolor{cinzaclaro}{HTML}{F2F5F7}
\definecolor{cinzaescuro}{HTML}{374151}
\hypersetup{colorlinks=true,linkcolor=azulpetroleo,urlcolor=azulpetroleo,citecolor=azulpetroleo}
\setlength{\parindent}{1.25cm}
\setlength{\parskip}{0.35em}
\onehalfspacing
\titleformat{\section}{\large\bfseries\color{azulpetroleo}}{}{0pt}{}
\titleformat{\subsection}{\normalsize\bfseries\color{azulpetroleo}}{}{0pt}{}
\pagestyle{fancy}
\fancyhf{}
\fancyhead[L]{\small Departamento de Eletroeletrônica — DEE}
\fancyhead[R]{\small IFMA — Campus São Luís/Monte Castelo}
\fancyfoot[C]{\small Página \thepage\ de \pageref{LastPage}}
\renewcommand{\headrulewidth}{0.4pt}
\renewcommand{\footrulewidth}{0.4pt}
\newcommand{\campo}[1]{\textcolor{ambar}{\textbf{[#1]}}}
\newcommand{\instituicao}{Instituto Federal de Educação, Ciência e Tecnologia do Maranhão}
\newcommand{\campus}{Campus São Luís/Monte Castelo}
\newcommand{\departamento}{Departamento de Eletroeletrônica — DEE}
\newcommand{\cidade}{São Luís — MA}
\newcommand{\cabecalhoInstitucional}{%
\begin{center}
{\Large\bfseries\color{azulpetroleo}\instituicao}\\[2mm]
{\large\bfseries\campus}\\[1mm]
{\normalsize\departamento}
\end{center}
\vspace{2mm}
\hrule
\vspace{4mm}
}
\newtcolorbox{destaque}{colback=cinzaclaro,colframe=azulpetroleo,boxrule=0.6pt,arc=2mm}

\begin{document}
\cabecalhoInstitucional
\begin{center}
{\Large\bfseries\color{azulpetroleo} Checklist do Procedimento de Doação e Parceria}
\end{center}

\section{Fase 1 — Planejamento interno}
\begin{itemize}[leftmargin=*]
\item[$\square$] Levantar necessidades dos laboratórios;
\item[$\square$] Consolidar equipamentos, quantidades e justificativas;
\item[$\square$] Identificar cursos, projetos e público beneficiado;
\item[$\square$] Preparar Projeto de Modernização Tecnológica;
\item[$\square$] Anexar planilha de empresas prioritárias.
\end{itemize}

\section{Fase 2 — Processo administrativo}
\begin{itemize}[leftmargin=*]
\item[$\square$] Abrir processo no SUAP;
\item[$\square$] Inserir solicitação formal do DEE;
\item[$\square$] Anexar projeto-mãe e plano de equipamentos;
\item[$\square$] Encaminhar à Direção-Geral;
\item[$\square$] Solicitar orientação/ciência dos setores administrativos e patrimoniais competentes;
\item[$\square$] Registrar autorização ou despacho para prospecção externa.
\end{itemize}

\section{Fase 3 — Prospecção das empresas}
\begin{itemize}[leftmargin=*]
\item[$\square$] Personalizar ofício para cada empresa;
\item[$\square$] Adequar os itens solicitados ao portfólio da empresa;
\item[$\square$] Identificar setor ou contato institucional correto;
\item[$\square$] Enviar ofício e projeto resumido;
\item[$\square$] Registrar data, responsável e retorno;
\item[$\square$] Agendar reunião técnica quando houver interesse.
\end{itemize}

\section{Fase 4 — Manifestação de interesse}
\begin{itemize}[leftmargin=*]
\item[$\square$] Obter manifestação formal da empresa;
\item[$\square$] Identificar bens, quantidades e valores;
\item[$\square$] Verificar estado, localização e titularidade dos bens;
\item[$\square$] Registrar eventuais encargos ou condicionantes;
\item[$\square$] Anexar documentação da empresa ao processo;
\item[$\square$] Submeter aos setores competentes do IFMA.
\end{itemize}

\section{Fase 5 — Formalização}
\begin{itemize}[leftmargin=*]
\item[$\square$] Verificar legislação e normas vigentes aplicáveis;
\item[$\square$] Verificar, quando cabível, procedimentos federais de manifestação de interesse/doação;
\item[$\square$] Realizar análise administrativa e patrimonial;
\item[$\square$] Realizar análise jurídica quando exigida;
\item[$\square$] Definir instrumento de formalização adequado;
\item[$\square$] Obter assinaturas das autoridades competentes;
\item[$\square$] Cumprir requisitos de publicidade/transparência aplicáveis.
\end{itemize}

\section{Fase 6 — Recebimento e patrimônio}
\begin{itemize}[leftmargin=*]
\item[$\square$] Conferir fisicamente os equipamentos;
\item[$\square$] Registrar modelos, números de série e quantidades;
\item[$\square$] Testar funcionamento quando aplicável;
\item[$\square$] Emitir relatório de recebimento;
\item[$\square$] Encaminhar ao patrimônio para registro/tombamento, quando cabível;
\item[$\square$] Registrar a destinação ao laboratório responsável;
\item[$\square$] Arquivar termos, notas, laudos e demais documentos no processo.
\end{itemize}

\section{Fase 7 — Pós-doação/parceria}
\begin{itemize}[leftmargin=*]
\item[$\square$] Instalar e colocar os equipamentos em operação;
\item[$\square$] Registrar treinamento/capacitação recebida;
\item[$\square$] Produzir registro fotográfico institucional;
\item[$\square$] Acompanhar uso em ensino, pesquisa e extensão;
\item[$\square$] Elaborar relatório de impacto quando pertinente;
\item[$\square$] Manter relacionamento institucional com a empresa parceira.
\end{itemize}

\begin{destaque}
\textbf{Atenção:} a prospecção pode ser conduzida pelo DEE, mas a aceitação formal da doação, assinatura do instrumento, registro patrimonial e demais atos devem observar a competência das autoridades e setores responsáveis do IFMA.
\end{destaque}
\end{document}
